\subsection{Projective dimension of a module}

DEFINITION 1. --- Let M be an A-module. One calls projective dimension of M, and denotes by dpA(M), the lower bound in Z of the lengths of projective resolutions of M (X, p.~48).

We therefore have \(\mathrm{dp}_{\mathbf{A}}(0) = -\infty\) , \(\mathrm{dp}_{\mathbf{A}}(\mathbf{M}) \geqslant 0\) if \(\mathbf{M} \neq 0\) . For \(\mathbf{M}\) to be projective, it is necessary and sufficient that \(\mathrm{dp}_{\mathbf{A}}(\mathbf{M}) \leqslant 0\) .

Lemma 1. --- If dp \(_{A}\) (M) \textless{} n \textless{} +∞, we have Ext \(_{A}^{n}\) (M, N) = 0 for every A-module N and Tor \(_{n}^{A}\) (P, M) = 0 for every right A-module P.

This follows immediately from the fact that M has a projective resolution of length \textless{} n and of X, p.~100, th. 1.

PROPOSITION 1. --- Let M be an A-module and n an integer ≥ 0. The following conditions are equivalent:

\begin{enumerate}
\def\labelenumi{(\roman{enumi})}
\item
  \(\mathrm{dp}_{\mathbf{A}}(\mathbf{M}) \leqslant n\) (i.e. (def. 1), M has a projective resolution of length \(\leqslant n\) )
\item
  \(\operatorname{Ext}_{\mathbf{A}}^{r}(\mathbf{M},\mathbf{N}) = 0\) for every A-module N and every integer \(r > n\) ;
\item
  \(\operatorname{Ext}_{\mathbf{A}}^{n+1}(\mathbf{M}, \mathbf{N}) = 0\) for every A-module \(\mathbf{N}\) ;
\item
  for every exact sequence
\end{enumerate}

\[
0 \rightarrow \mathrm{K} \rightarrow \mathrm{P} _ {n - 1} \rightarrow \dots \rightarrow \mathrm{P} _ {0} \rightarrow \mathrm{M} \rightarrow 0
\]

where the \(P_{i}\) are projective, K is projective.

\begin{enumerate}
\def\labelenumi{(\roman{enumi})}
\item
  \(\Rightarrow\) (ii): this follows from Lemma 1.
\item
  \(\Rightarrow\) (iii): this is trivial.
\item
  \(\Rightarrow\) (iv): in the situation of (iv), one has for every A-module N an isomorphism of \(\operatorname{Ext}_{\mathbf{A}}^{1}(\mathbf{K},\mathbf{N})\) over \(\operatorname{Ext}_{\mathbf{A}}^{n + 1}(\mathbf{M},\mathbf{N})(\mathbf{X},\mathbf{p}.128,\operatorname{cor.}4)\) ; if (iii) is satisfied, one has
\end{enumerate}

\[
\operatorname{Ext} _ {\mathbf {A}} ^ {1} (\mathbf {K}, \mathbf {N}) = 0
\]

for all N and K is projective (X, p.~93, prop. 10).

\begin{enumerate}
\def\labelenumi{(\roman{enumi})}
\setcounter{enumi}{3}
\tightlist
\item
  \(\Rightarrow\) (i): consider the exact sequence (X, p.~50)
\end{enumerate}

\[
0 \to Z _ {n - 1} (\mathrm{M}) \to \mathrm{L} _ {n - 1} (\mathrm{M}) \to \dots \to \mathrm{L} _ {0} (\mathrm{M}) \to \mathrm{M} \to 0.
\]

If (iv) is satisfied, \(Z_{n-1}(M)\) is projective and M has a projective resolution of length \(\leqslant n\) .

COROLLARY 1. --- Let \((\mathbf{M}_i)_{i\in \mathbb{E}}\) be a family of A-modules. We have

\[
\mathrm{dp} _ {\mathbf {A}} \left(\bigoplus_ {i \in \mathbf {E}} \mathbf {M} _ {i}\right) = \sup _ {i \in \mathbf {E}} \mathrm{dp} _ {\mathbf {A}} (\mathbf {M} _ {i}).
\]

This follows from the equivalence of conditions (i) and (iii) of Prop. 1, and from Prop. 7 of X, p.~89.

In the following statement, we agree that \(\pm\infty + 1 = \pm\infty - 1 = \pm\infty\) .

COROLLARY 2. --- Let

\[
0 \rightarrow \mathrm{M} ^ {\prime} \rightarrow \mathrm{M} \rightarrow \mathrm{M} ^ {\prime \prime} \rightarrow 0
\]

be an exact sequence of A-modules.

\begin{enumerate}
\def\labelenumi{\alph{enumi})}
\tightlist
\item
  We have
\end{enumerate}

\[
\mathrm{dp} _ {\mathbf {A}} (\mathbf {M}) \leqslant \sup \left(\mathrm{dp} _ {\mathbf {A}} (\mathbf {M} ^ {\prime}), \mathrm{dp} _ {\mathbf {A}} (\mathbf {M} ^ {\prime \prime})\right).
\]

Equality holds as soon as \(\mathrm{dp}_{\mathbf{A}}(\mathbf{M}^{\prime \prime})\neq \mathrm{dp}_{\mathbf{A}}(\mathbf{M}^{\prime}) + 1\)

\begin{enumerate}
\def\labelenumi{\alph{enumi})}
\setcounter{enumi}{1}
\tightlist
\item
  We have
\end{enumerate}

\[
\mathrm{dp} _ {\mathbf {A}} \left(\mathrm{M} ^ {\prime \prime}\right) \leqslant \sup \left(\mathrm{dp} _ {\mathbf {A}} (\mathrm{M}), \mathrm{dp} _ {\mathbf {A}} \left(\mathrm{M} ^ {\prime}\right) + 1\right).
\]

Equality holds as soon as \(\mathrm{dp}_{\mathbf{A}}(\mathbf{M}) \neq \mathrm{dp}_{\mathbf{A}}(\mathbf{M}')\) .

\[
\mathrm{dp} _ {\mathbf {A}} \left(\mathbf {M} ^ {\prime}\right) \leqslant \sup \left(\mathrm{dp} _ {\mathbf {A}} (\mathbf {M}), \mathrm{dp} _ {\mathbf {A}} \left(\mathbf {M} ^ {\prime \prime}\right) - 1\right)
\]

Equality holds as soon as \(\mathrm{dp}_{\mathbf{A}}(\mathbf{M}) \neq \mathrm{dp}_{\mathbf{A}}(\mathbf{M}^{\prime \prime})\) .

Let us prove, for example, part a), the proofs of parts b) and c) being analogous.

\[
\geqslant 0.
\]

\[
\begin{array}{r l}\operatorname{Ext} _ {\mathbf {A}} ^ {n + 1} (\mathrm{M} ^ {\prime \prime}, \mathrm{N})&\rightarrow \operatorname{Ext} _ {\mathbf {A}} ^ {n + 1} (\mathrm{M}, \mathrm{N}) \rightarrow \operatorname{Ext} _ {\mathbf {A}} ^ {n + 1} (\mathrm{M} ^ {\prime}, \mathrm{N}) \rightarrow \operatorname{Ext} _ {\mathbf {A}} ^ {n + 2} (\mathrm{M} ^ {\prime \prime}, \mathrm{N})\\&\rightarrow \operatorname{Ext} _ {\mathbf {A}} ^ {n + 2} (\mathrm{M}, \mathrm{N}).\end{array}
\]

If \(\mathrm{dp}_{\mathbf{A}}(\mathbf{M}')\) , \(\mathrm{dp}_{\mathbf{A}}(\mathbf{M}'')\leqslant n\) , then \(\operatorname{Ext}_{\mathbf{A}}^{n+1}(\mathbf{M}',\mathbf{N})=0\) and \(\operatorname{Ext}_{\mathbf{A}}^{n+1}(\mathbf{M}''\) , \(\mathbf{N})=0\) (prop. 1) therefore \(\operatorname{Ext}_{\mathbf{A}}^{n+1}(\mathbf{M},\mathbf{N})=0\) and \(\mathrm{dp}_{\mathbf{A}}(\mathbf{M})\leqslant n\) (prop. 1), so that

\[
\mathrm{dp} _ {\mathbf {A}} (\mathbf {M}) \leqslant \sup \left(\mathrm{dp} _ {\mathbf {A}} (\mathbf {M} ^ {\prime}), \mathrm{dp} _ {\mathbf {A}} (\mathbf {M} ^ {\prime \prime})\right).
\]

If \(\mathrm{dp}_{\mathbf{A}}(\mathbf{M}) < \sup (\mathrm{dp}_{\mathbf{A}}(\mathbf{M}^{\prime}),\mathrm{dp}_{\mathbf{A}}(\mathbf{M}^{\prime \prime}))\) then necessarily \(\mathrm{dp}_{\mathbf{A}}(\mathbf{M}) < + \infty\) . For every \(n > \mathrm{dp}_{\mathbf{A}}(\mathbf{M})\) and for every A-module N, we have

\[
\operatorname{Ext} _ {\mathbf {A}} ^ {n + 1} \left(\mathrm{M} ^ {\prime}, \mathrm{N}\right) \neq 0 \Leftrightarrow \operatorname{Ext} _ {\mathbf {A}} ^ {n + 2} \left(\mathrm{M} ^ {\prime \prime}, \mathrm{N}\right) \neq 0,
\]

by the preceding exact sequence; by prop. 1, this immediately implies \(\mathrm{dp}_{\mathbf{A}}(\mathbf{M}^{\prime \prime}) = \mathrm{dp}_{\mathbf{A}}(\mathbf{M}^{\prime}) + 1\) , since one of the quantities \(\mathrm{dp}_{\mathbf{A}}(\mathbf{M}^{\prime})\) , \(\mathrm{dp}_{\mathbf{A}}(\mathbf{M}^{\prime \prime})\) is \(> \mathrm{dp}_{\mathbf{A}}(\mathbf{M})\) .

Example. --- Let \(a\) be an element of \(\mathbf{A}\) which is neither invertible nor a right zero divisor. Then \(\mathrm{dp}_{\mathbf{A}}(\mathbf{A}/\mathbf{A}a)=1\) .

Indeed, by the exact sequence \(0 \to \mathbf{A}_s \xrightarrow{\varphi} \mathbf{A}_s \to \mathbf{A} / \mathbf{A}a \to 0\) , where \(\varphi(x) = xa\) , we have \(\mathrm{dp}_{\mathbf{A}}(\mathbf{A} / \mathbf{A}a) \leqslant 1\) . If \(\mathrm{dp}_{\mathbf{A}}(\mathbf{A} / \mathbf{A}a) < 1\) , then \(\mathbf{A} / \mathbf{A}a\) is projective, and there exists an A-linear map \(\psi : \mathbf{A}_s \to \mathbf{A}_s\) such that \(\psi \circ \varphi = 1d\) ; this implies

\[
1 = \psi (\varphi (1)) = \psi (a) = a. \psi (1)
\]

and a is invertible.

PROPOSITION 2. --- Suppose A is left noetherian. Let M be a finitely generated A-module and n an integer ≥ 0. The following conditions are equivalent:

\begin{enumerate}
\def\labelenumi{(\roman{enumi})}
\tightlist
\item
  \(\mathrm{dp}_{\mathbf{A}}(\mathbf{M}) \leqslant n\) .
\end{enumerate}

(i bis) M has a projective resolution P of length ≤ n such that \(P_{i}\) is a finitely generated A-module for each i.

\begin{enumerate}
\def\labelenumi{(\roman{enumi})}
\setcounter{enumi}{1}
\item
  \(\operatorname{Ext}_{\mathbf{A}}^{r}(\mathbf{M},\mathbf{N}) = 0\) for every finitely generated A-module N and every integer \(r > n\) .
\item
  \(\operatorname{Ext}_{\mathbf{A}}^{n+1}(\mathbf{M}, \mathbf{N}) = 0\) for every finitely generated A-module \(\mathbf{N}\) .
\item
  \(\operatorname{Tor}_r^{\mathbf{A}}(\mathbf{P},\mathbf{M}) = 0\) for every right A-module P and every \(r > n\) .
\item
  \(\operatorname{Tor}_{n+1}^{\mathbf{A}}(\mathbf{A}/\mathfrak{a},\mathbf{M}) = 0\) for every finitely generated right ideal \(\mathfrak{a}\) of \(\mathbf{A}\) .
\end{enumerate}

(i bis) \(\Rightarrow\) (i): this is trivial.

\begin{enumerate}
\def\labelenumi{(\roman{enumi})}
\item
  \(\Rightarrow\) (ii): this follows from Lemma 1.
\item
  \(\Rightarrow\) (iii): this is trivial.
\item
  \(\Rightarrow\) (i): by (iii) and X, p.~107, prop. 5, we have \(\operatorname{Ext}_{\mathbf{A}}^{n+1}(\mathbf{M}, \mathbf{N}) = 0\) for every A-module N, whence (i) by prop. 1.
\item
  \(\Rightarrow\) (iv): this follows from lemma 1.
\item
  \(\Rightarrow\) (v): this is trivial.
\item
  \(\Rightarrow\) (i bis): let \((\mathbf{L}, d)\) be a free resolution of \(\mathbf{M}\) such that \(\mathbf{L}_r\) be finitely generated for every \(r\) (X, p.~53, prop. 6). Set \(\mathbf{K} = \mathbf{Z}_{n-1}(\mathbf{L})\) ; then \(\mathbf{K}\) is finitely generated as a submodule of \(\mathbf{L}_{n-1}\) and we have an exact sequence
\end{enumerate}

\[
0 \rightarrow \mathrm{K} \rightarrow \mathrm{L} _ {n - 1} \rightarrow \mathrm{L} _ {n - 2} \rightarrow \dots \rightarrow \mathrm{L} _ {1} \rightarrow \mathrm{L} _ {0} \rightarrow \mathrm{M} \rightarrow 0.\tag{1}
\]

By (v) and X, p.~131, cor. 3, we have \(\operatorname{Tor}_{1}^{\mathbf{A}}(\mathbf{A}/\mathfrak{a},\mathbf{K}) = 0\) for every finitely generated right ideal \(\mathfrak{a}\) of A. According to th. 2 of X, p.~74, the A-module K is flat; since it is finitely generated, hence finitely presented (X, p.~10, prop. 5), it is projective (X, p.~13, cor.), so (1) is a projective resolution of M.

COROLLARY. --- Suppose A is left noetherian and let \(\mathcal{C}_{0}\) (resp. \(\mathcal{C}\) ) be the set of classes of finitely generated projective A-modules (resp. of finitely generated A-modules of finite projective dimension). Then the homomorphism of Grothendieck groups \(\mathrm{K}(\mathcal{C}_{0}) \to \mathrm{K}(\mathcal{C})\) is bijective.

This follows from X, p.~58, th. 1 (note that \(\mathcal{C}_0\) and \(\mathcal{C}\) are left exact by cor. 2).

